\subsection{交换幺半群的分式幺半群}

在本号中，e 将表示幺半群的单位元，而 \(x^{*}\) 表示 E 的可逆元素 x 的逆。

设 E 是一个交换幺半群，S 是 E 的一个子集，S' 是由 S 生成的 E 的子幺半群。

引理1. 在 \(\mathbf{E} \times \mathbf{S}'\) 中，由以下方式定义的关系 \(\mathbf{R}_{\xi}^{x}, y_{\xi}^{x}\) ：

“存在 \(a, b \in E\) 和 \(p, q, s \in S'\) ，使得 \(x = (a, p)\) , \(y = (b, q)\) 且 \(aqs = bps\) ''

是一个与乘积幺半群 \(\mathbf{E} \times \mathbf{S}'\) 上的律相容的等价关系.

显然 R 是自反的且对称的。设 \(x = (a, p)\) , \(y = (b, q)\) 和 \(z = (c, r)\) 为 \(E \times S'\) 的元素，使得 \(R\xi x, y\xi\) 和 \(R\xi y, z\xi\) 成立。则存在 \(S'\) 的两个元素 s 和 t，使得

\[
a q s = b p s, \quad b r t = c q t,
\]

由此可得

\[
a r (s t q) = b p s r t = c p (s t q)
\]

且因此 \(\mathbf{R}\{x, z\}\) 成立，因为 \(stq\) 属于 \(\mathbf{S}'\) 。关系 \(\mathbf{R}\) 因此是传递的。

此外，令 \(x = (a, p), y = (b, q), x' = (a', p')\) ，且 \(y' = (b', q')\) 为 \(E \times S'\) 的元素，使得 \(R\xi x, y\xi\) 和 \(R\xi x', y\xi\) 成立。存在 \(s\) 和 \(s'\) （在 \(S'\) 中），使得

\[
a q s = b p s, \quad a ^ {\prime} q ^ {\prime} s ^ {\prime} = b ^ {\prime} p ^ {\prime} s ^ {\prime},
\]

由此可得 \((aa')(qq')(ss') = (bb')(pp')(ss')\) ，因此 \(\mathbf{R}\{\chi x', yy'\}\) 成立（对于 \(ss' \in S'\) ）。等价关系 \(\mathbf{R}\) 因此与 \(\mathbf{E} \times \mathbf{S}'\) 上的合成律相容.

商广群 \((\mathbf{E} \times \mathbf{S}')/\mathbf{R}\) 是一个交换幺半群。

定义7. 设 E 为交换幺半群，S 为 E 的子集，S' 为 E 中由 S 生成的子幺半群。商幺半群 (E × S')/R（其中等价关系 R 如引理1所述）记为 E \(_{S}\) ，称为 E 的与 S 相伴的分式幺半群†（或分母在 S 中的分式幺半群）。

对于 \(a \in \mathbf{E}\) 和 \(p \in \mathbf{S}'\) ， \((a, p)\) 模 \(\mathbf{R}\) 的类通常记作 \(a/p\) ，称为分子为 \(a\) 、分母为 \(p\) 的分式。于是根据定义 \((a/p)\) . \((a'/p') = aa'/pp'\) 。分式 \(a/p\) 和 \(a'/p'\) 相等，当且仅当存在 \(s\) （在 \(\mathbf{S}'\) 中）使得 \(spa' = sp'a\) ；若如此，则存在 \(\sigma\) 和 \(\sigma'\) （在 \(\mathbf{S}'\) 中）使得 \(a\sigma = a'\sigma'\) 和 \(p\sigma = p'\sigma'\) 。特别地， \(a/p = sa/sp\) （对于 \(a \in \mathbf{E}\) 和 \(s, p\) （在 \(\mathbf{S}'\) 中））。 \(E_S\) 的单位元是分式 \(e/e\) .

设 \(a / e = \varepsilon(a)\) （对所有 \(a \in E\) ）。上述表明 \(\varepsilon\) 是 \(E\) 到 \(E_S\) 的一个同态，称为典范同态。对所有 \(p \in S'\) , \((p / e) \cdot (e / p) = e / e\) ，因此 \(e / p\) 是 \(\varepsilon(p) = p / e\) 的逆；于是 \(\varepsilon(S')\) 的每个元素都是可逆的。于是 \(a / p = (a / e) \cdot (e / p)\) ，由此

\[
a / p = \varepsilon (a) \varepsilon (p) ^ {*}\tag{1}
\]

对于 \(a \in \mathbf{E}\) 和 \(p \in \mathbf{S}'\) ，幺半群 \(\mathbf{E}_{\mathbf{S}}\) 因此由 \(\varepsilon(\mathbf{E}) \cup \varepsilon(\mathbf{S})^*\) 生成.

命题6. 记号与定义7相同， \(\varepsilon\) 表示 E 到 \(\mathbf{E}_{\mathrm{S}}\) 的典范同态。

\begin{enumerate}
\def\labelenumi{(\roman{enumi})}
\item
  设 \(a\) 和 \(b\) （属于 \(\mathbf{E}\) ）；要使 \(\varepsilon(a) = \varepsilon(b)\) ，其充要条件是存在 \(s \in \mathbf{S}'\) 使得 \(sa = sb\) .
\item
  \(\varepsilon\) 为单射的充要条件是 S 中的每个元素都是可消去的。
\item
  \(\varepsilon\) 为双射的充要条件是 \(S\) 的每个元素都可逆。
\end{enumerate}

断言(i)是显然的，并且表明 \(\varepsilon\) 是单射，当且仅当 \(S'\) 的每个元素都是可消去的；但由于 \(E\) 的可消去元素所成之集是 \(E\) 的一个子幺半群（第2号，命题2），这等同于说 \(S\) 的每个元素都是可消去的。

若 \(\varepsilon\) 是双射的，则 S 的每个元素都是可逆的，因为 \(\varepsilon(S)\) 由 \(E_{S}\) 的可逆元素组成。反之，假设 S 的每个元素都可逆；那么 \(S'\) 的每个元素都可逆（第3号，命题4的推论2），因此可消去。于是由(ii)知 \(\varepsilon\) 是单射，且由(1)有 \(a/p = \varepsilon(a.p^{*})\) ，因此 \(\varepsilon\) 是满射。

定理1. 设 E 为交换幺半群，S 为 E 的子集， \(E_{S}\) 为与 S 相伴的分式幺半群， \(\varepsilon: E \to E_{S}\) 为典范同态。进一步设 f 为 E 到幺半群 F（不一定是交换的）的一个同态，使得 \(f(S)\) 的每个元素在 F 中可逆。则存在且仅存在一个同态 \(\bar{f}\) （从 \(E_{S}\) 到 F），使得 \(f = \bar{f} \circ \varepsilon\) .

若 \(\bar{f}\) 是 \(\mathbf{E}_{\mathbf{S}}\) 到 \(\mathbf{F}\) 的一个同态，使得 \(f = \bar{f} \circ \varepsilon\) ，则

\[
\bar {f} (a | p) = \bar {f} (\varepsilon (a) \varepsilon (p) ^ {*}) = \bar {f} (\varepsilon (a)) \bar {f} (\varepsilon (p)) ^ {*} = f (a) f (p) ^ {*}
\]

对所有 \(a \in \mathbf{E}\) 和 \(p \in \mathbf{S}'\) 成立，由此得到 \(\bar{f}\) 的唯一性.

设 \(g\) 是 \(\mathbf{E} \times \mathbf{S}'\) 到 \(\mathbf{F}\) 的映射，定义为 \(g(a, p) = f(a).f(p)^*\) 。我们证明 \(g\) 是 \(\mathbf{E} \times \mathbf{S}'\) 到 \(\mathbf{F}\) 的一个同态。首先，

\[
g (e, e) = f (e) f (e) ^ {*} = e.
\]

设 \((a, p)\) 和 \((a', p')\) 为 \(\mathbf{E} \times \mathbf{S}'\) 的两个元素；由于 \(a'\) 和 \(p\) 在 \(\mathbf{E}\) 中交换， \(f(a')\) 和 \(f(p)\) 在 \(\mathbf{F}\) 中交换，因此 \(f(a')f(p)^* = f(p)^*f(a')\) （依据第3号，命题5）。此外 \(f(p p')^* = f(p' p)^* = (f(p')f(p))^* = f(p)^*f(p')^*\) （依据第3号，命题4），因此

\[
\begin{array}{l} g (a a ^ {\prime}, p p ^ {\prime}) = f (a a ^ {\prime}) f (p p ^ {\prime}) ^ {*} = f (a) f (a ^ {\prime}) f (p) ^ {*} f (p ^ {\prime}) ^ {*} = f (a) f (p) ^ {*} f (a ^ {\prime}) f (p ^ {\prime}) ^ {*} \\ = g (a, p) g (a ^ {\prime}, p ^ {\prime}). \end{array}
\]

我们证明 \(g\) 与等价关系 \(\mathbf{R}\) （在 \(\mathbf{E} \times \mathbf{S}'\) 上）相容：如果 \((a, p)\) 和 \((a', p')\) 模 \(\mathbf{R}\) 同余，则存在 \(s \in \mathbf{S}'\) 使得 \(spa' = sap'\) ，因此 \(f(s)f(p)f(a') = f(s)f(a)f(p')\) 。由于 \(f(s)\) 是可逆的，由此可得 \(f(p)f(a') = f(a)f(p')\) ，然后左乘 \(f(p)^*\) 且右乘 \(f(p')^*\)

\[
g (a ^ {\prime}, p ^ {\prime}) = f (a ^ {\prime}) f (p ^ {\prime}) ^ {*} = f (p) ^ {*} f (a) = f (a) f (p) ^ {*} = g (a, p).
\]

因此存在一个同态 \(\bar{f}\) （从 \(\mathbf{E}_{\mathbf{S}}\) 到 \(\mathbf{F}\) ），使得 \(\bar{f}(a/p) = g(a,p)\) ，于是 \(\bar{f}(\varepsilon(a)) = \bar{f}(a/e) = f(a)f(e)^* = f(a)\) 。因此 \(\bar{f} \circ \varepsilon = f\) .

推论. 设 E 和 F 是两个交换幺半群，S 和 T 分别是 E 和 F 的子集，f 是从 E 到 F 的同态，使得 \(f(\mathrm{S}) \subset \mathrm{T}\) ， \(\varepsilon: E \to E_{S}\) , \(\eta: F \to F_{T}\) 为典范同态。存在且仅存在一个同态 \(g: E_{S} \to F_{T}\) ，使得 \(g \circ \varepsilon = \eta \circ f\) .

同态 \(\eta \circ f\) （从 \(E\) 到 \(F_{T}\) ）将 \(S\) 的每个元素映射到 \(F_{T}\) 的一个可逆元素。

注。(1) 定理1也可以表述为： \((\mathbf{E}_{\mathbf{S}},\varepsilon)\) 是 \(\mathbf{E}\) 的泛映射问题的解，该问题相对于幺半群、幺半群同态以及 \(\mathbf{E}\) 到幺半群的同态（它们将 \(\mathbf{S}\) 的元素映为可逆元素）（集合论，IV，§ 3，第 1 号）。由此得出（同上），该问题的每一个其他解都以唯一方式同构于 \((\mathbf{E}_{\mathbf{S}},\varepsilon)\) .

\begin{enumerate}
\def\labelenumi{(\arabic{enumi})}
\setcounter{enumi}{1}
\tightlist
\item
  为使上述泛映射问题有解，不必假定幺半群 E 是交换的，这由集合论，IV，§3，第2号（参见习题17）可知。
\end{enumerate}

我们提及分式幺半群的两个重要特例。

\begin{enumerate}
\def\labelenumi{(\alph{enumi})}
\item
  设 \(\overline{\mathbf{E}} = \mathbf{E}_{\mathbb{E}}\) 。由于幺半群 \(\overline{\mathbf{E}}\) 由集合 \(\varepsilon(\mathbf{E}) \cup \varepsilon(\mathbf{E})^*\) （由可逆元素组成）生成， \(\overline{\mathbf{E}}\) 的每个元素都是可逆的（第3号，命题4的推论2）。换言之， \(\overline{\mathbf{E}}\) 是一个交换群。此外，根据定理1，每个同态 \(f\) （从 \(\mathbf{E}\) 到群 \(\mathbf{G}\) ）都可以唯一地分解为 \(f = \bar{f} \circ \varepsilon\) 的形式，其中 \(\bar{f}: \overline{\mathbf{E}} \to \mathbf{G}\) 是一个同态。 \(\overline{\mathbf{E}}\) 称为 \(\mathbf{E}\) 的分式群（在加法记号下称为 \(\mathbf{E}\) 的差群）。
\item
  设 \(\Phi = \mathrm{E}_{\Sigma}\) ，其中 \(\Sigma\) 由 \(\mathbf{E}\) 的可消去元素组成。根据命题6(ii)， \(\mathbf{E}\) 到 \(\Phi\) 的典范同态是单射的；将 \(\mathbf{E}\) 与它在 \(\Phi\) 中的像等同起来是有益的。于是 \(\mathbf{E}\) 是 \(\Phi\) 的一个子幺半群， \(\mathbf{E}\) 的每个可消去元在 \(\Phi\) 中都有逆元，且 \(\Phi\) 的每个元素都具有形式 \(a/p = a.p^{*}\) （其中 \(a \in \mathbf{E}\) 和 \(p \in \Sigma\) ）；于是 \(a/p = a'/p'\) 当且仅当 \(ap' = pa'\) 。容易看出， \(\Phi\) 的可逆元素是分式 \(a/p\) （其中 \(a\) 和 \(p\) 都可消去），且 \(p/a\) 是 \(a/p\) 的逆.
\end{enumerate}

现在设 S 是 E 的一组可消去元素，S' 是由 S 生成的 E 的子幺半群。若 \(a/p\) 和 \(a'/p'\) 是 \(\mathbf{E}_{\mathbf{S}}\) 的两个元素，则 \(a/p = a'/p'\) 当且仅当 \(ap' = pa'\) （因为 \(sap' = spa'\) 蕴含 \(ap' = pa'\) ，对所有 \(s \in \mathbf{S}'\) ）。 \(\mathbf{E}_{\mathbf{S}}\) 因此可以与 \(\Phi\) 的由 \(\mathbf{E} \cup \mathbf{S}^*\) 生成的子幺半群等同。

如果 E 的每个元素都是可消去的，则 \(\Phi = \bar{E}\) 且 E 是交换群 \(\Phi\) 的一个子幺半群。反之，若 E 同构于某个群的子幺半群，则 E 的每个元素都是可消去的。

